% ============================================================
%  Conclusion
% ============================================================
\renewcommand{\tablename}{Table}
\renewcommand{\figurename}{Figure}

\chapter*{Conclusion}
\addcontentsline{toc}{chapter}{Conclusion}

This study proposes and evaluates an Independent Multi-Agent Reinforcement
Learning (Independent MARL) system for traffic signal control, in which each
signalised intersection is managed by a separate agent using a shared policy.
Three algorithms --- PPO, DQN, and DDQN --- are combined with three reward
functions --- queue, pressure, and wait-clip --- and evaluated on two scenarios
extracted from the real-world TAPAS Cologne traffic dataset: Cologne-3
(3 signalised intersections) and Cologne-8 (8 signalised intersections).

On Cologne-8, all RL configurations outperform both the Fixed-time and Webster
baselines on mean queue length and throughput. The best configuration,
PPO-pressure, reduces mean queue by 88\,\% relative to Fixed-time and 96\,\%
relative to Webster. These results demonstrate that Independent MARL can
effectively react to local congestion at each intersection, overcoming
Webster's limitation of optimising each node independently without accounting
for inter-intersection flow interactions.

Regarding reward functions, pressure is the best choice when combined with PPO
or DQN, while queue is the safest and most consistent choice across all three
algorithms. Wait-clip performs well with PPO but causes severe instability with
DDQN due to its sparse and delayed reward signal. In terms of practical
deployability, queue is the most feasible reward function, as the number of
stopped vehicles at an intersection can be estimated directly from cameras or
inductive loop sensors --- infrastructure already widely available at signalised
junctions. Pressure requires sensors on both incoming and outgoing lanes,
which is feasible but demands more complete sensing infrastructure. Wait-clip
requires tracking the cumulative waiting time of individual vehicles,
which is considerably more complex and difficult to deploy under current
traffic infrastructure conditions.

On Cologne-3, the RL configurations do not demonstrate consistent advantage.
The training curves show no clear convergence over the full 500{,}000 steps,
and evaluation results exhibit high standard deviations across most
configurations. One possible contributing factor is that fewer intersections
means fewer samples are accumulated per rollout for the shared policy within
the same timestep budget, making convergence more difficult. This suggests that
Independent MARL achieves its clearest advantage when the number of signalised
intersections is large enough for the shared policy to accumulate sufficient
and diverse training signal.

Several limitations of the present study should be acknowledged. All training
and evaluation are conducted within the SUMO simulation environment and have
not been deployed on real traffic signal hardware; the simulation does not
capture sensor noise, device latency, or real-world driver behaviour. The
experiments are also limited to two small networks with at most 8 intersections,
and the scalability of the approach to larger networks with hundreds of nodes
has not been verified. Finally, the unstable results on Cologne-3 indicate
that the current approach is not well suited to small networks with few
intersections, where the learning signal is sparser within the same timestep
budget.

Several directions for future work follow from these limitations. Testing
scalability on larger networks is necessary to assess the practical viability
of the method; extending the comparison to other RL algorithms such as
Soft Actor-Critic (SAC)~\cite{haarnoja2018sac} and Twin Delayed Deep
Deterministic Policy Gradient (TD3)~\cite{fujimoto2018td3} may provide further
insight into which algorithm families are best suited to large-scale traffic
signal control. For Cologne-3, a direct improvement would be to increase the
training timestep budget to accumulate sufficient learning signal, or to adopt
MARL architectures with more explicit coordination such as
MAPPO~\cite{yu2022mappo} or QMIX~\cite{rashid2018qmix}, in which agents can
share information during training rather than acting in full independence.

% ------------------------------------------------------------
%  BibTeX entries to add to your .bib file
% ------------------------------------------------------------
%
% @inproceedings{haarnoja2018sac,
%   author    = {Tuomas Haarnoja and Aurick Zhou and Pieter Abbeel
%                and Sergey Levine},
%   title     = {Soft Actor-Critic: Off-Policy Maximum Entropy Deep
%                Reinforcement Learning with a Stochastic Actor},
%   booktitle = {Proceedings of the 35th International Conference on
%                Machine Learning (ICML)},
%   year      = {2018}
% }
%
% @inproceedings{fujimoto2018td3,
%   author    = {Scott Fujimoto and Herke van Hoof and David Meger},
%   title     = {Addressing Function Approximation Error in
%                Actor-Critic Methods},
%   booktitle = {Proceedings of the 35th International Conference on
%                Machine Learning (ICML)},
%   year      = {2018}
% }
%
% @inproceedings{yu2022mappo,
%   author    = {Chao Yu and Akash Velu and Eugene Vinitsky and
%                Jiaxuan Gao and Yu Wang and Alexandre Bayen and Yi Wu},
%   title     = {The Surprising Effectiveness of {PPO} in Cooperative
%                Multi-Agent Games},
%   booktitle = {Advances in Neural Information Processing Systems
%                (NeurIPS)},
%   year      = {2022}
% }
%
% @inproceedings{rashid2018qmix,
%   author    = {Tabish Rashid and Mikayel Samvelyan and
%                Christian Schroeder de Witt and Gregory Farquhar and
%                Jakob Foerster and Shimon Whiteson},
%   title     = {{QMIX}: Monotonic Value Function Factorisation for
%                Deep Multi-Agent Reinforcement Learning},
%   booktitle = {Proceedings of the 35th International Conference on
%                Machine Learning (ICML)},
%   year      = {2018}
% }